\section{Discussion and Limitations}
\label{sec:limits}

\paragraph{What the results say.}
Under attainable caps, most models cannot reach first place under any of the six rules, and the rule determines how many can: the share at 10 points ranges from 3\% to 29\% on OLLB and from 18\% to 51\% on HELM.
Cost comparisons across rules are meaningful only where the same targets are feasible under all of them; there, the minimum and median rules need the fewest points, Borda the most tasks, and the mean rule's task count depends on the cap.
The ordering of rules differs between OLLB and HELM and between caps (Figures~\ref{fig:feas} and~\ref{fig:cap}), so the supported statement is conditional: it holds per cap scenario and scale, not as a single ranking of rules.
The closed forms identify why: the minimum rule needs every task below the best rival's level to reach it, the median rule a majority, and the mean rule only a total, which any cap that does not bind can supply from few tasks.

\paragraph{Limitations.}
\emph{The estimand is a score counterfactual, not a training cost.}
It assumes that any score in $[x_{it},c_t]$ is reachable independently on each task, with no spillover (gains on one task need not transfer, and may reduce other scores), no rival response and no equilibrium, and it says nothing about the data, compute or contamination that would deliver the gain.
\emph{Caps are assumptions.}
Frontier and empirical caps are not evidence that such scores are attainable; the empirical caps come from observed fine-tune-minus-base pairs (survivorship bias, one quantile for all targets), and on HELM they are a proxy borrowed from another benchmark.
The conclusions on feasibility are therefore as good as the cap scenario; we show eight and rank nothing across them.
\emph{Scale and size.}
OLLB has $n=6$ tasks on unrelated scales, so median, minimum and geometric-mean results are tied to the official normalisation (Appendix Figure~\ref{fig:scale}) and the median and Condorcet results are tie-sensitive; HELM is one benchmark with 57 same-unit subjects and 79 models, and describes a snapshot, not a population of leaderboards.
\emph{Cohorts and intervals.}
Borda and Condorcet on OLLB are computed on a 136-target stratified sample (stratum-weighted); cohorts at attainable caps are small (3 to 91 targets); the Condorcet cost is an upper bound for any Condorcet-consistent rule; $K^\ast$ has few levels at $n=6$.
The task bootstrap and the evaluation-noise simulations use 10 targets and are sensitivity analyses conditional on the model and caps; they do not validate them.
\emph{Margin.}
Strict victory over a continuous score has an infimum, so we fix $\varepsilon=10^{-3}$ and report sensitivity to it (Table~\ref{tab:eps}).
\emph{Measurement and data.}
Scores are single-run leaderboard numbers; one OLLB model per name is kept; the licences of the score data are not declared on the sources we could verify.
We use no evidence about any model's training data, and nothing here accuses any listed model.

\paragraph{Broader impact.}
The study documents how little targeted improvement some rules require and how others are out of reach, which can inform more robust aggregation and also show where leaderboards are easiest to game.
It uses public score tables only, with no human subjects and no model training.

\paragraph{Conclusion.}
First place under a multi-task rule is not a single price: it is a feasibility question first and a cost question only among targets for which first place is attainable.
On the two matrices studied, the aggregation rule changes how many models can win and what it costs, and neither ordering is stable across datasets and caps.
